\section{The CAO-DF Decision Rule}
\label{sec:framework}

CAO-DF turns the cost model into a decision. It uses no machine learning: every decision can
be traced back to a measured saving, a measured one-time cost and the value of $N$.

\textbf{Inputs and output.} The rule takes the features of a new workload $w$, the expected
number of runs $N$, and a table of measured \emph{reference} workloads. For every reference
workload and action, the table stores two numbers: the saving per run as a share of a static
run, $s = \Delta/\Cstatic$, and the one-time cost in static runs, $r = \Copt/\Cstatic$. The
output is \emph{keep static}, \emph{tune} (A1) or \emph{compact} (A2).

\textbf{Estimating an unseen workload.} We describe each workload by three features: the
logarithm of the number of orders, the logarithm of the number of input files, and the share
of rows that reach the shuffle. Each feature is scaled by its range over the reference
workloads. For each action, CAO-DF finds the \emph{nearest} reference workload (Euclidean
distance; equal distances are averaged) and copies its $s$ and $r$. We chose this method
because it is the simplest one that can be explained in one sentence; threshold rules and
interpolation were the alternatives.

\textbf{Decision.} In units of one static run, the predicted total cost is $N$ for keeping the
static configuration and $r + N(1 - s)$ for an action (Fig.~\ref{fig:algo}). An action is only
considered if its estimated saving is at least a noise floor $\delta$; the cheapest option
wins, and ties go to \emph{keep static}.

\begin{figure}[t]
\centering
\fbox{\begin{minipage}{0.94\columnwidth}
\small\raggedright
\textbf{Input:} workload $w$, runs $N$, reference table $T$, floor $\delta$\\
\textbf{Output:} $a^{*} \in \{\textit{static}, \mathrm{A1}, \mathrm{A2}\}$
\begin{algorithmic}[1]
  \STATE $TC(\textit{static}) \leftarrow N$
  \FOR{$a \in \{\mathrm{A1}, \mathrm{A2}\}$}
    \STATE $(s_a, r_a) \leftarrow$ values of the nearest workload to $w$ in $T$ for $a$
    \STATE $TC(a) \leftarrow r_a + N(1 - s_a)$ \textbf{if} $s_a \ge \delta$ \textbf{else} $\infty$
  \ENDFOR
  \RETURN $\arg\min_{a} TC(a)$, ties $\rightarrow$ \textit{static}
\end{algorithmic}
\end{minipage}}
\caption{The CAO-DF decision rule (costs in units of one static run).}
\label{fig:algo}
\end{figure}

\textbf{Why a noise floor.} Without it, the rule would act on random variation. Compacting a
single-file input saves nothing, yet its measured ``saving'' can be slightly positive by chance
and give a finite break-even point (Section~\ref{sec:results}). We set $\delta = 10\,\%$ of a
static run. In our data, the spread (interquartile range) of the static runs was 0.2--16\,\%
of their median, every real saving was at least 23\,\%, and every noise ``saving'' was at most
2\,\%, so 10\,\% separates the two clearly.

\textbf{Baselines.} We compare CAO-DF with three alternatives: \emph{static} (always keep the
default), \emph{performance-only} (always pick the action with the largest estimated saving,
ignoring its one-time cost), and an \emph{oracle} that knows the measured costs in hindsight.
